\documentclass[12pt]{article}
\usepackage[utf8]{inputenc}
\usepackage[T1]{fontenc}
\usepackage{hyperref}
\title{Fali, South Language Resources for K-12 Curriculum Development}
\author{Compiled by Claude Sonnet 5.5 (claude-sonnet-5-5) for Andr\'as Kornai}
\date{}
\begin{document}
\maketitle

\section{Introduction}
South Fali (ISO 639-3: fal) is spoken in northern Cameroon. Wikipedia gives about 20,000 speakers; Joshua Project lists a larger figure of 68,000 for the people group. Varieties include Fali-Tinguelin, Kaang and Bele. Its Adamawa affiliation is disputed. It is distinct from North Fali. The searches found no dictionary, primer or school programme for South Fali.

\section{Primary Educational Resources}
None found.

\section{Linguistic and Reference Materials}
\begin{itemize}
\item Wikipedia article on the Fali languages of Cameroon: \url{https://en.wikipedia.org/wiki/Fali_languages_(Cameroon)}
\item Raija Kramer, \emph{Die Sprache der Fali in Nordkamerun} (2014, a German functional grammar with phonology and tone; publisher page): \url{https://www.koeppe.de/titel_die-sprache-der-fali-in-nordkamerun}
\item General Alphabet of Cameroon Languages: \url{https://en.wikipedia.org/wiki/General_Alphabet_of_Cameroon_Languages}
\end{itemize}

\section{Religious and Community Resources}
\begin{itemize}
\item Scripture Earth listing for ISO fal: \url{https://www.scriptureearth.org/00eng.php?iso=fal}
\item Joshua Project language page for Fali, South: \url{https://joshuaproject.net/index.php/languages/fal}
\end{itemize}

\section{Digital and Audio Resources}
None beyond the Scripture Earth listing above, whose contents were not verified.

\section{Curriculum Development Context}
There is a reference grammar (in German) but no pedagogical materials were found. Cameroon's General Alphabet is the likely base for an orthography, but no confirmed South Fali orthography was found. Development would begin with the grammar, the Scripture Earth holdings, and community consultation.
\end{document}